\scanpage{41}
\deleted{\{ $\unsure{\mathcal C|\mathcal E}$. catégorie \{ des $k$-modules $1/2$ simpliciaux \deleted{\unsure{tels}}\inserted{$L_*$} qui
sont flasques et str. $\partial$-parfaits, i.e. qui satisfont
deux \deleted{aux} conditions équivalentes

a) $L_*$ est \deleted{\unsure{équivalent}}\inserted{isomorphe} à une somme directe finie
de modules $1/2$ simpliciaux
$\deleted{\unsure{\Phi^i_*=\Lambda^i\Phi_*|\Phi(i)_*=\Lambda^i\Phi_*}\ (i\geq1)}\inserted{C^i_*=\Lambda^{i+1}\Phi_{*k}\ (i\geq0)}$
}
% slanted
\margin{\deleted{On suppose tt module proj. de type fini sur $k$ est libre p. ex. $k$ principal}}
\deleted{b) Les $L_n$ sont libres de type fini, de rang constant
sur $\operatorname{Spec} k$, la fonction $n\mapsto\operatorname{rg}_k L_n$ est polynomiale
pour $n$ grand, $L_*$ est flasque i.e. $\pi_i(L_*)=0$ pour $i\geq0$.

$\mathcal A\subset\mathcal C$ la sous-catégorie pleine formée des
$\deleted{\unsure{\Phi^i_*}}=C^{i-1}_*$

\uline{NB} $\mathcal C$ est additive (et karoubienne), $\mathcal A$ n'est pas
stable par sommes finies, mais hérite une structure
$k$-linéaire. On a
\[
 \operatorname{Hom}(\deleted{\unsure{\Phi}}C^i_*,\deleted{\unsure{\Phi}}C^j_*)=
 \begin{cases}
 0&\text{si }j\ne i,i+1\\
 k&\text{si }j=i,\ \text{base }\mathrm{id}_{\unsure{\Phi(i)|C^i|C^i_*}}\\
 k&\text{si }j=i+1,\ \text{base }d_i=T_*\wedge
 \end{cases}
\]
Notons que
\deleted{$\unsure{\Phi_*}\ \unsure{\Phi^i_*|C^i_*}=C^{i+1}_*$ \unsure{inversement} $C^i_*=\unsure{\Phi^{i+1}_*}$ $(i\geq0)$}

$C^\bullet_*=(C^i_*,T_*\wedge)$ est un complexe de cochaînes dans
$\mathcal A$ ou dans $\mathcal C$ ou dans $\widehat\Delta_{k\text{-}\mathrm{lin}}$ -- il jouera
le rôle du complexe de cochaînes \uline{universel}.}

\deleted{\unsure{En effet}}

\uline{Prop} 1. Soit $\mathcal B$ une catégorie \deleted{additive}\inserted{\deleted{\unsure{lin.}} $k$-additive}. Alors
\diagram{\begin{tikzcd}[ampersand replacement=\&,column sep=2em,row sep=1.4em]
\underline{\operatorname{Hom}}_{\deleted{\unsure{\mathrm{add}|k\text{-}\mathrm{lin}}}\unsure{k\text{-}\mathrm{lin}}}(\unsure{\mathcal K|\mathcal C},\mathcal B)
 \arrow[r,"{\substack{\mathrm{rest}\\\simeq}}"] \&
 \underline{\operatorname{Hom}}_{k\text{-}\mathrm{lin}}(\mathcal A,\mathcal B)
 \arrow[d,"{\simeq}"',"{F\mapsto F(C^\bullet_*)}"] \\
\& {\mathcal B^\bullet_{\unsure{k}}=\text{cat. des \inserted{complexes de} cochaînes de }\mathcal B}
\end{tikzcd}}

\deleted{\uline{\unsure{Question|Corollaire}} Soit $K^\bullet\in\operatorname{Ob}\mathcal B^*$ un complexe de cochaînes de $\mathcal B$}

\uline{Exemple 1} \inserted{faisons $\mathcal B=\mathrm{Ab}_k$.} Soit $L_*\in\operatorname{Ob}\mathcal C$, considérons le
foncteur cov. qu'il représente
\[
 \mathcal K\longrightarrow\mathrm{Ab}_k:\qquad M_*\longmapsto\operatorname{Hom}_{\mathcal C}(L_*,M_*)
\]
on a ainsi un foncteur pl. fid. (sous cat. des fonct. repr. --)
\diagram{\begin{tikzcd}[ampersand replacement=\&,row sep=0.9em,column sep=1em]
\mathcal K^\circ \arrow[r,hook] \& \underline{\operatorname{Hom}}_{k\text{-}\mathrm{lin}}(\mathcal K,\mathrm{Ab}_k)
 \arrow[dr,"{\simeq}"] \& \& \\
\& \& \underline{\operatorname{Hom}}_{k\text{-}\mathrm{lin}}(\mathcal A,\mathrm{Ab}_k) \arrow[r,"{\simeq}"] \& \mathrm{Ab}^{\bullet\deleted{\unsure{*}}}_{k\deleted{\unsure{*}}}
\end{tikzcd}}
% archive sheet number: 40
